\label{sec:hierarchy}
\subsection{Eldan's cumulant hierarchy is the law of total cumulance in time}
Let $\mu$ be a compactly supported probability measure on $\R^d$. Run Eldan's normalized localization: $\mu_t\propto
e^{\langle c_t,x\rangle-\frac12\langle Q_tx,x\rangle}\mu$ with $dc=A_t^{-1}a_tdt+A_t^{-1/2}dB_t$ and $dQ=A_t^{-1}dt$. Write $\kappa_m(t)$ for the
cumulant tensors of $\mu_t$.
\begin{proposition}[Bizeul--Klartag--Lehec, Lemma 4.3; known]\label{prop:bkl}
$d\kappa_m=\kappa_{m+1}(\cdot,\dots,\cdot,A_t^{-1/2}dB_t)-(m\kappa_m+L_m)\,dt$, with $L_m$ as in Section~\ref{sec:readings}. In particular:
\begin{itemize}[leftmargin=*]
\item $L_3=0$, so $\E\kappa_3$ only decays, at rate $3$.
\item The fourth-order term is
\[
L_4(h_1,\dots,h_4)=\sum_{j=2}^4\big\langle\kappa_3(h_1,h_j,\cdot),A^{-1}\kappa_3(h_{\{2,3,4\}\setminus\{j\}},\cdot)\big\rangle .
\]
For instance, $L_4(u,u,u,u)=3\langle\kappa_3(u,u,\cdot),A^{-1}\kappa_3(u,u,\cdot)\rangle$.
\end{itemize}
\end{proposition}
\begin{proposition}[the drift is total cumulance; proved]\label{prop:totalcum}
$\mu_{dt}$ is the posterior of $X\sim\mu$ given one infinitesimal Gaussian observation. Brillinger's law of total cumulance over
set partitions $\pi$,
\[
\kappa_m(X)=\sum_\pi\kappa_{|\pi|}\big(\kappa_{B}(X\,|\,\mathcal G)\big)_{B\in\pi},
\]
applied to that observation, reproduces the drift $-(m\kappa_m+L_m)$ exactly:
\begin{itemize}[leftmargin=*]
\item the one-block partition gives $\E\kappa_m(dt)$;
\item two-block partitions with a singleton give $m\kappa_m\,dt$;
\item two-block partitions with both blocks of size $\ge2$ give $L_m\,dt$;
\item partitions into three or more blocks contribute $o(dt)$.
\end{itemize}
\end{proposition}
\begin{proof}
Conditionally on the observation, $\kappa_k(dt)=\kappa_k+\kappa_{k+1}(\cdot,A^{-1/2}dB)+O(dt)$. The covariance over $dB$ of the fluctuations of
$\kappa_{|B_1|}$ and $\kappa_{|B_2|}$ is $\langle\kappa_{|B_1|+1},A^{-1}\kappa_{|B_2|+1}\rangle dt$. A singleton block $B_1=\{i\}$ has $\kappa_2=A$, and
$\langle A h_i,A^{-1}\kappa_m(h_{\rm rest},\cdot)\rangle=\kappa_m(h)$. Joint cumulants of three or more $O(\sqrt{dt})$ Gaussian-linear fluctuations
vanish at order $dt$.
\end{proof}
So $L_m$ is the \emph{between-path} part of the $m$-th cumulant: what conditioning on the localization path moves out of the
within-path cumulants. It is quadratic in lower cumulants, and the metric is the inverse covariance, because the observation noise is
$A^{-1/2}dB$. The identity is checked on a $7$-point measure in $\R^3$ by Gauss--Hermite averaging over one step (Richardson
in $dt$).
\begin{itemize}[leftmargin=*]
\item The drift of $\kappa_3$ matches $-3\kappa_3$ to $1.5\times10^{-7}$.
\item The drift of $\kappa_4$ matches $-(4\kappa_4+L_4)$ to $2.8\times10^{-7}$; dropping $L_4$ errs by $21\%$.
\item $L_4(u^4)$ equals $3\langle\kappa_3(u,u,\cdot),A^{-1}\kappa_3(u,u,\cdot)\rangle$ to all digits.
\end{itemize}

\subsection{The Schur hub is the between-tier fourth cumulant}
\begin{proposition}[static form; proved]\label{prop:static}
For any $\sigma$-algebra $\mathcal G$ and scalar $X$,
\begin{multline*}
\kappa_4(X)=\E\kappa_4(X|\mathcal G)+4\Cov\!\big(\E[X|\mathcal G],\kappa_3(X|\mathcal G)\big)+3\Var\!\big(\Var(X|\mathcal G)\big)\\
+6\,\kappa\big(\E[X|\mathcal G],\E[X|\mathcal G],\Var(X|\mathcal G)\big)+\kappa_4\big(\E[X|\mathcal G]\big).
\end{multline*}
If $\mathcal G\supseteq\sigma(Y)$ for a vector $Y$ with covariance $C_Y$, then the projection inequality gives
\[
3\Var(\Var(X|\mathcal G))\ \ge\ 3\,\Cov(\Var(X|\mathcal G),Y)\,C_Y^{-1}\,\Cov(Y,\Var(X|\mathcal G)).
\]
Equality holds when the conditional variance is affine in $Y$.
\end{proposition}
Read in the two-tier language of stage 12:
\begin{itemize}[leftmargin=*]
\item the upper tier is the conditioning ($\mathcal G$, the localization centre, i.e.\ the field the next layer sees);
\item the lower tier is the conditional law;
\item the term $3\Var(\Var(X|\mathcal G))$ is the \emph{between-tier} fourth cumulant, and its Schur projection is the
$\kappa_3A^{-1}\kappa_3$ form of $L_4$.
\end{itemize}
The hierarchy's term has the same shape. For centred $X$, $\kappa_3(u,u,\cdot)=\Cov((u\cdot X)^2,X)$, so
\[
L_4(u^{\otimes4})=3\,\Cov((u\cdot X)^2,X)\,A^{-1}\,\Cov(X,(u\cdot X)^2)=3\Var\big(\mathrm{proj}_{\mathrm{span}X}(u\cdot X)^2\big)
\]
under $\mu_t$, averaged over paths: three times the variance of the linear regression of the square on the field. v56's path class is a
regression of the same kind. Note XLI writes it as
\[
P4_i=12\,[YC^{-1}Y^\top]_{ii}=3\Var\big(\mathrm{proj}_{\mathrm{span}\,z_l}\,\mathsf C_i\big),\qquad 2Y_{ij}=\Cov(\mathsf C_i,z_{l,j}),
\]
where $\mathsf C_i$ is the cross term between neuron $i$'s first and second chaos. Using the cross term instead of the whole square removes two
within-tier pieces: the square of the first chaos, which the Gaussian closure already handles, and the square of the second chaos, the
$4$-cycle.

In the chaos-2 model $z=L\cdot x+\frac12(x^\top Hx-\tr H)$, the cross term projects onto the input's linear span as $2(HL)\cdot x$, so
$3\Var$ of it is $12|HL|^2$, the path class of $\kappa_4=12|HL|^2+3\tr H^4$. Under localization this projection is the part of the
conditional variance that moves linearly with the centre. Indeed $\Var(z|\mathcal F_t)=|L+Hm_t|^2/(1+t)+\tr H^2/2(1+t)^2$, whose
linear fluctuation is $2L^\top Hm_t/(1+t)$. So the path class is the between-tier part of the fourth cumulant, read through the field.

\begin{corollary}[what it explains in v56; interpretation, with note XLI's data]
\begin{enumerate}[leftmargin=*]
\item \textbf{The metric is $C^{-1}$, not $\diag(C)^{-1}$.} The observation noise of localization is $A^{-1/2}dB$.
Measured: the hub with the diagonal metric is $30$--$60$ times too large.
\item \textbf{The cumulants must be conditional (quenched).} The drift involves $\E[\kappa_3(t)A_t^{-1}\kappa_3(t)]$ over paths, not
$\E\kappa_3\,A^{-1}\E\kappa_3$. The annealed slice $D21=\kappa(z_i,z_i,\cdot)$ mixes within-tier and between-tier pieces. Measured: the
symmetrized-D21 control removes $0$--$6\%$, the source-resolved hub $64$--$81\%$.
\item \textbf{The weight is fixed.} The three pairings of $L_4(e_i^{\otimes4})$ are the $3$ in $3\Var$. Measured: free-running prefers the
derived weight $1$, and renormalization leaves it at $1.007$.
\item \textbf{Each half alone is adverse.} Proposition~\ref{prop:static} has five terms, and the third-chaos star is the
$4\Cov(\E,\kappa_3)$-type partner. Measured: either piece alone is $+79\%$ to $+131\%$; together they give $-21\%$.
\end{enumerate}
\end{corollary}

\subsection{Predictions: the other pairings from the same solve}
The hierarchy gives every slice of $L_4$, with coefficients:
\begin{gather*}
L_4(e_i,e_i,e_j,e_j)=\langle D_i,A^{-1}D_j\rangle+2\langle\kappa_3(e_i,e_j,\cdot),A^{-1}\kappa_3(e_i,e_j,\cdot)\rangle,\\
L_4(e_i,e_i,e_i,e_j)=3\langle D_i,A^{-1}\kappa_3(e_i,e_j,\cdot)\rangle,
\end{gather*}
with $D_i=\kappa_3(e_i,e_i,\cdot)$. In v56's dressed form the first term of the $(2,2)$ slice is note XLI's ``cheap path part''
$4[YM^{-1}Y^\top]_{ij}$ (one product from the existing solve). The second term and the $(3,1)$ slice need the off-diagonal third cumulant
contracted once: $\sum_k\kappa_3(i,j,k)\,y_{ik}$ with $y_i=C^{-1}D_i$ already formed. On the carried CP sources this is one hub-type
contraction per source-layer, about the cost of the existing D21 hub.
\begin{conjecture}[testable on note XLI's dumps]
The source-resolved $(3,1)$ pairing
\[
3\sum_k\kappa_3(i,j,k)\,[C^{-1}Y]_{ik}
\]
removes a substantial part of the chain's K31 error. That
error is $74\%$ rms; about $40\%$ of the (3,1) slice is ``history'', not a function of the pair state. The pairing uses the full source
tensor, not the pair state, so it is exactly the kind of term the flex theorem allows to carry new information.
\end{conjecture}
\begin{remark}[$\kappa_3$ has no Schur term]
$L_3=0$: under localization the expected third cumulant only decays. A per-layer closure of $\kappa_3$ therefore has no
between-tier correction to add. Its errors are transport errors. This matches note XXXIX: the chain's D21 error is the flat
incoherent sector of the transport (the product-gate defect), right to $2\%$ in its coherent part.
\end{remark}
